% Plain Dharma — XeLaTeX preamble for the screen PDF.
% Loaded by pandoc via --include-in-header. Pandoc's default book template
% supplies \maketitle, \tableofcontents, hyperref, microtype, etc. — this file
% only adds the project-specific tuning: fonts, palette, chapter typography.

% Local Garamond Libre fonts (same files served by the website).
% Family name has no space so the `*` wildcard expands to the actual filename
% (e.g. `*-Regular` → `GaramondLibre-Regular.otf`).
\setmainfont{GaramondLibre}[
  Path=__FONTS_DIR__/,
  Extension=.otf,
  UprightFont=*-Regular,
  ItalicFont=*-Italic,
  BoldFont=*-Bold,
  BoldItalicFont=*-BoldItalic,
  Ligatures=TeX,
  Numbers=OldStyle
]

% Brand palette — matches the site's --color-* tokens.
\definecolor{cream}{HTML}{F5EFE0}
\definecolor{ink}{HTML}{2A2520}
\definecolor{muted}{HTML}{6B5F50}
\definecolor{saffron}{HTML}{C7651C}
\definecolor{linkcolor}{HTML}{8A4F18}

% Cream page background — matches the cover artwork and the website. Without
% this the cover (cream) → inner pages (white) transition is jarring.
\pagecolor{cream}

% Body color and slightly looser line spacing for reading on screen.
\color{ink}
\linespread{1.18}

% Chapter heading: no "Chapter N" prefix, smaller and quieter than LaTeX default.
\usepackage{titlesec}
\titleformat{\chapter}[display]
  {\normalfont\Huge\bfseries\color{ink}}
  {}
  {0pt}
  {\Huge}
\titlespacing*{\chapter}{0pt}{40pt}{30pt}

% Section heading: also calmer than default.
\titleformat{\section}
  {\normalfont\Large\bfseries\color{ink}}
  {}
  {0pt}
  {}

% Blockquote (used by the per-chapter drops): italic, indented, muted color.
\renewenvironment{quote}{%
  \begin{list}{}{%
    \leftmargin=1.5em\rightmargin=1.5em\itemindent=0pt%
  }\item\itshape\color{muted}}{%
  \end{list}}

% Cap images at 70% of the text width. Pandoc wraps standalone-paragraph
% images in its own centering — overriding \includegraphics here breaks
% inline-context placements (e.g. images inside quote blocks), so we only
% touch the default size.
\usepackage{graphicx}
\setkeys{Gin}{width=0.7\linewidth,keepaspectratio}

% Garamond Libre is missing some Pali diacritics (ṇ U+1E47, ṭ U+1E6D, etc).
% Fall back to STIX Two Text — a serif with full Unicode coverage that pairs
% visually with Garamond — for just those characters. Avoids the unsightly
% missing-glyph blank in chapter subtitles like "Ādittapariyāya Sutta".
\usepackage{newunicodechar}
\newfontfamily\palifallback{STIX Two Text}
\newunicodechar{ṇ}{{\palifallback ṇ}}
\newunicodechar{ṭ}{{\palifallback ṭ}}

% Tighter list spacing (pandoc's default is roomy).
\usepackage{enumitem}
\setlist{itemsep=0.15em,topsep=0.3em}

% Don't break paragraphs with hanging indents — use space between paragraphs.
\setlength{\parindent}{0pt}
\setlength{\parskip}{0.6em}

% Hyphenation hints — chapter titles shouldn't break the brand or compound
% words at awkward places ("Mind-fulness" looked bad in the first render).
\hyphenation{Mindfulness Loving-Kindness Dharma Buddha Buddhist}

% Hyperlink colors are wired via pandoc's `-V colorlinks/linkcolor/urlcolor`
% CLI flags — pandoc loads hyperref *after* this header, so calling
% \hypersetup directly here would fire before the package exists.
% Use \AtBeginDocument to defer until hyperref is loaded.
\AtBeginDocument{%
  \hypersetup{
    pdfsubject={Six Foundational Buddhist Teachings in Plain Modern English},
    pdfkeywords={Buddhism, suttas, dharma, meditation, mindfulness},
    pdfcreator={pandoc + XeLaTeX}
  }%
}
